%% Согласно ГОСТ Р 7.0.11-2011:
%% 5.3.3 В заключении диссертации излагают итоги выполненного
% исследования, рекомендации, перспективы дальнейшей разработки темы.
%% 9.2.3 В заключении автореферата диссертации излагают итоги данного
% исследования, рекомендации и перспективы дальнейшей разработки темы.
\begin{enumerate}
  \item На основе анализа архитектурных особенностей
    многопроцессорных систем и существующих моделей памяти (SC, TSO,
    XC, RC, RVWMO) предложена математическая модель представления
    многопоточной программы в виде частично упорядоченного множества
    операций с памятью. Такое представление позволяет свести задачу
    анализа консистентности памяти к задаче построения и анализа
    линейных порядков, сохраняя при этом все ограничения,
    накладываемые моделью памяти.
  \item Численные исследования показали, что прямое построение всех
    линейных порядков частично упорядоченного множества сталкивается
    с комбинаторным взрывом: в худшем случае количество порядков
    растёт как $n!$, а средняя асимптотическая сложность алгоритма
    составляет $O(2^n)$. Это подтвердило необходимость разработки
    специализированных алгоритмов и структур данных для сокращения
    пространства поиска.
  \item Математическое моделирование показало, что учёт классов
    эквивалентности перестановок (независимости и зависимости
    операций) позволяет существенно сократить количество
    анализируемых линейных порядков. Разработаны алгоритмы
    топологической сортировки с учётом независимости
    (алгоритмы~\ref{alg:topo-indep},~\ref{alg:topo-indep-lookup}) и с
    учётом зависимостей (алгоритм~\ref{alg:topo-dep}). Получены
    оценки сложности, показывающие снижение вычислительных затрат до
    $O(n^2)$ в предельном случае высокой независимости операций.
  \item Для выполнения поставленных задач был создан комплекс методов
    и структур данных, включающий:
    \begin{itemize}
      \item структуру данных ``граф состояний'' (State Graph) для
        компактного и полного представления всех возможных
        результатов исполнения многопоточной программы, позволяющую
        восстанавливать линейные порядки без их явного хранения;
      \item алгоритмы верификации соответствия аппаратуры заявленной
        модели памяти как на основе линейных порядков, так и на
        основе анализа графа состояний;
      \item метод кластеризации графа состояний и метрику насыщения,
        позволяющие декомпозировать задачу верификации на независимые
        подзадачи и различать эффекты, обусловленные
        переупорядочиванием зависимых операций и моделью памяти.
    \end{itemize}
\end{enumerate}

Практическая значимость подтверждена экспериментальной апробацией на
тестовых сценариях, моделирующих литмус-тесты и нагрузки с большим
числом конкурирующих обращений к памяти. Показано, что предложенный
подход позволяет масштабировать верификацию на программы с десятками
тысяч операций, что выходит за рамки применимости классических подходов.

Перспективы дальнейшей разработки темы. Проведённое исследование
открывает несколько направлений для дальнейшей работы. Одно из них
связано с адаптацией предложенного подхода к гетерогенным
вычислительным архитектурам: современные системы всё чаще объединяют
ядра общего назначения с ускорителями (GPU, FPGA, специализированные
тензорные блоки), обладающими собственной иерархией памяти и
собственными моделями согласованности. Это порождает задачу построения
композиционной модели памяти, в которой граф состояний строится для
каждой подсистемы отдельно, а затем объединяется через точки
синхронизации, что позволило бы распространить предложенный метод на
гетерогенные платформы без потери масштабируемости. Другое направление
связано с развитием эвристик анализа: метрика насыщения, предложенная
в работе, позволяет останавливать исполнение программы при достижении
полного покрытия кластеров состояний, и дальнейшее развитие этого
направления состоит в разработке адаптивных стратегий генерации
тестовых сценариев, которые на основе текущего распределения
маркированных кластеров целенаправленно увеличивают покрытие редких
состояний, в том числе потенциально содержащих ошибки реализации
подсистемы памяти.
